\documentclass[10pt,a4paper]{article}
\usepackage[T1]{fontenc}
\usepackage[utf8]{inputenc}
\usepackage{lmodern,microtype}
\usepackage[margin=0.72in,headheight=14pt,headsep=15pt,footskip=24pt]{geometry}
\usepackage{amsmath,amssymb,booktabs,tabularx,array,graphicx,xcolor,enumitem,fancyhdr,titlesec,caption}
\usepackage[hyphens]{url}
\usepackage[colorlinks=true,linkcolor=navy,urlcolor=navy,citecolor=navy]{hyperref}
\definecolor{navy}{HTML}{193D59}
\definecolor{muted}{HTML}{56616B}
\pagestyle{fancy}\fancyhf{}
\fancyhead[L]{\small\textcolor{muted}{Small-model alignment: SFT and DPO}}
\fancyhead[R]{\small\textcolor{muted}{Zain Abdullah Ahmad | 27100215}}
\fancyfoot[C]{\small\thepage}
\renewcommand{\headrulewidth}{0.3pt}
\setlength{\parindent}{0pt}\setlength{\parskip}{5pt}
\setlength{\tabcolsep}{5pt}\renewcommand{\arraystretch}{1.17}
\titleformat{\section}{\large\bfseries\color{navy}}{\thesection}{0.55em}{}
\titleformat{\subsection}{\normalsize\bfseries\color{navy}}{\thesubsection}{0.5em}{}
\titlespacing*{\section}{0pt}{9pt}{5pt}\titlespacing*{\subsection}{0pt}{7pt}{3pt}
\captionsetup{font=small,labelfont=bf,skip=4pt,hypcap=false}
\newcolumntype{Y}{>{\raggedright\arraybackslash}X}
\newcommand{\code}[1]{\texttt{\small\detokenize{#1}}}
\newcommand{\fig}[3]{\begin{center}\includegraphics[width=\linewidth]{figures/#1.png}\captionof{figure}{#2}\label{#3}\end{center}}
\newcommand{\case}[2]{\vspace{4pt}\textbf{#1}\quad\textcolor{muted}{\small #2}\par}
\hypersetup{pdftitle={Instruction Tuning and Direct Preference Optimization on Small Language Models},pdfauthor={Zain Abdullah Ahmad},pdfsubject={Programming Assignment 2: experimental report}}
\begin{document}
\thispagestyle{plain}
{\LARGE\bfseries\color{navy} Instruction Tuning and Direct Preference Optimization on Small Language Models\par}
\vspace{6pt}
{\large Zain Abdullah Ahmad\quad\textcolor{muted}{27100215}\par}
{\small Generative AI: Programming Assignment 2\hfill 16 October 2026\par}
\vspace{5pt}\hrule\vspace{6pt}
\textbf{Abstract.} This experiment examines how preference training changes a small language model after supervised fine-tuning (SFT). A custom SFT loop trained Qwen2.5-0.5B on 30,000 instruction--response examples, after which two independent Direct Preference Optimization (DPO) runs used 15,000 HH-RLHF or UltraFeedback pairs. Both began from the same SFT checkpoint and used a frozen copy as their reference. In blinded comparisons, HH DPO won 190 times against 78 SFT wins and 32 ties, giving a 70.90\% win rate excluding ties. UltraFeedback DPO performed worse than SFT, with 103 wins, 172 losses and 25 ties (37.45\%). A separate persona LoRA experiment on SmolLM2 also performed worse than its base model, winning 20.51\% of non-tied joint persona--quality comparisons and 11.11\% of general-quality comparisons. The saved responses help explain these differences. HH DPO stopped more reliably, UltraFeedback DPO often exhausted the output budget, and persona LoRA repeatedly generated loops. Across these runs, improved validation preference metrics were insufficient evidence of better generated responses.

\section{Repository, models, and data}
The code and completed runs are available at \href{https://github.com/zainahmadlums/genai-pa2}{github.com/zainahmadlums/genai-pa2}, with the reported evidence preserved at commit \href{https://github.com/zainahmadlums/genai-pa2/tree/20cb609f71e9d032b7b5ceee0a4644123c7f5ed9}{\texttt{20cb609}}. The main experiments use \code{Qwen/Qwen2.5-0.5B} with custom SFT and DPO training loops. The separate persona experiment uses \code{HuggingFaceTB/SmolLM2-360M-Instruct} with TRL and PEFT. The assignment starter supplied the persona preference data and synthetic HH evaluation set.

\begin{center}\small
\begin{tabularx}{\linewidth}{@{}lYrr@{}}\toprule
Track & Dataset / pinned revision prefix & Train & Validation\\\midrule
SFT & \href{https://huggingface.co/datasets/HuggingFaceH4/ultrachat_200k}{HuggingFaceH4/ultrachat\_200k}, \texttt{train\_sft}; \texttt{8049631c405a} & 30,000 & 1,000\\
HH DPO & \href{https://huggingface.co/datasets/Anthropic/hh-rlhf}{Anthropic/hh-rlhf}; \texttt{09be8c5bbc57} & 15,000 & 1,000\\
UF DPO & \href{https://huggingface.co/datasets/argilla/ultrafeedback-binarized-preferences-cleaned}{argilla/ultrafeedback-binarized-preferences-cleaned}; \texttt{770076f077c4} & 15,000 & 1,000\\
Persona & Provided \texttt{ramsay\_style\_preferences.json} & 200 & 20\\\bottomrule
\end{tabularx}
\captionof{table}{Actual training and validation subsets. UF denotes UltraFeedback. Full revisions, selection counts and file hashes are preserved in the provenance records.}
\end{center}

Each training track uses single-turn examples. UltraChat contributes the first user--assistant exchange, while HH pairs must contain exactly one human turn and one assistant turn with a shared prompt on both sides. To give the HH run a substantial safety component, selection begins from a balanced harmless/helpful pool. The resulting training set contains 7,423 harmless and 7,577 helpful pairs, with helpful examples drawn from the base, online and rejection-sampled subsets. UltraFeedback pairs are retained when the chosen response has a higher rating than the rejected response.

The same length and overlap filters apply to all three Qwen tracks. Selection removes duplicate normalized instructions and known evaluation instructions, requires at least two response tokens, and limits formatted prompts to 256 tokens. The complete prompt, response and EOS must fit within 512 tokens. These filters preserve complete response targets, although they also favor shorter examples.

Seed 42 controls data selection, splitting, training, generation setup and balanced A/B placement. Evaluation uses 300 provided synthetic HH prompts, divided equally among ten categories, and 40 fixed persona prompts covering cooking, coding, general help and safety. For general evaluation, 300 prompts were sampled without replacement from the 805 AlpacaEval items at revision \texttt{2edc6fad8be6}, using seed 42. All evaluation files are committed under \code{eval_sets/}. The AlpacaEval results apply to this fixed subset and judging procedure; they are not official leaderboard scores.

\newpage
\section{Experimental configuration}
SFT trains the model to predict instruction responses. Two independent DPO runs then train from the selected SFT checkpoint, using HH or UltraFeedback preferences and the same frozen SFT reference. DPO uses summed response log-probabilities with $\beta=0.1$~\cite{dpo}. Table~\ref{tab:config} gives the final settings.

\begin{center}\small
\begin{tabularx}{\linewidth}{@{}lYYYY@{}}\toprule
Run & Strategy / reference & Optimizer, learning rate, schedule & Batch / precision / length & Epochs; steps; $\beta$\\\midrule
SFT & Full tuning; no reference & AdamW, $2\!\times\!10^{-5}$; cosine; 3\% warmup; WD 0.1 & $2\times16=32$; FP16 autocast; 512 & 1; 938/938; --\\
HH DPO & Full tuning; frozen SFT best, step 938 & AdamW, $10^{-6}$; constant; no warmup; WD 0 & $1\times32=32$; FP16 autocast; 512 & 1; 469/466; 0.1\\
UF DPO & Full tuning; frozen SFT best, step 938 & AdamW, $10^{-6}$; constant; no warmup; WD 0 & $1\times32=32$; FP16 autocast; 512 & 1; 469/466; 0.1\\
Persona DPO & LoRA: $r=16$, $\alpha=32$, dropout 0.05; frozen base, adapter disabled & AdamW (Trainer default), $10^{-4}$; cosine; 10\% warmup; WD 0 & $4\times2=8$; FP16 autocast; 512 & 3; 75 steps; 0.1\\\bottomrule
\end{tabularx}
\captionof{table}{Batch is micro-batch $\times$ gradient accumulation. Qwen steps are attempted/applied updates. All runs use seed 42 and clipping at 1.0. Evaluation uses greedy decoding, 256 new tokens, repetition penalty 1.0, no-repeat n-gram 0 and no added length penalty. Temperature/top-$p$ are logged as 1.0 but inactive. The prompt budget is 1,024 tokens; Qwen uses the template below and persona uses the base chat template.}
\label{tab:config}
\end{center}

All runs keep trainable weights in FP32 and use FP16 autocast with gradient scaling. The Qwen DPO policy and frozen reference occupy separate GPUs, both with FP32 weights. SFT selects minimum validation loss; Qwen DPO selects maximum reward accuracy, breaking ties by reward margin. Validation occurs initially, every 100 attempted steps and at completion. Persona selects minimum validation loss, evaluated every ten steps.

The persona adapter targets attention projections \code{q_proj}, \code{k_proj}, \code{v_proj}, \code{o_proj} and MLP projections \code{gate_proj}, \code{up_proj}, \code{down_proj}~\cite{peft,trl}. Its 8,683,520 trainable parameters comprise 2.34\% of the 370,504,640 total. Rank 16 limits adaptation capacity, while dropout provides light regularization. Attention and MLP targets allow changes across contextual and lexical processing. The $10^{-4}$ learning rate was chosen to adapt the newly initialized low-rank parameters while the base stayed frozen. An effective batch of eight gives 25 updates per epoch on 200 pairs. Three epochs allow repeated exposure to this small dataset, with validation selecting among checkpoints. These were initial choices, without a tuning sweep.

Every Qwen stage uses the following template without a system prompt:
\begin{center}\small\ttfamily
\#\#\# Instruction:\textbackslash n\{instruction\}\textbackslash n\textbackslash n\#\#\# Response:\textbackslash n
\end{center}
SFT loss is averaged over response tokens, including EOS. Persona training and generation use SmolLM2's tokenizer chat template with its default system message, ``You are a helpful AI assistant named SmolLM, trained by Hugging Face'', and no additional persona instruction.

\section{Parameter choices and deviations}
The Qwen learning rates, $\beta$, training batches, accumulation, epoch counts and clipping follow the guide. Resource settings retain its one-policy-GPU arrangement and FP16 autocast. The guide permits an FP16 reference; this run uses FP32 reference weights under the same autocast as the policy to match arithmetic, yielding tied initial rewards.

The remaining choices concern data and evaluation. The unspecified SFT prompt cap becomes 256 tokens, reserving response space. Permitted truncation becomes complete-example filtering at 512 tokens, preserving targets but selecting shorter examples. The unspecified HH mixture becomes a balanced candidate pool, yielding 49.5\% harmless training pairs to emphasize safety. The guide does not require judge hints; this evaluation includes the supplied expected-behavior text to make its safety criteria explicit. Their individual quality effects were not isolated.

\newpage
\section{Training behavior}
\fig{sft}{SFT loss and held-out perplexity. Blue curves show logged training-window means; orange points are validation values. No additional smoothing is applied. The dashed line marks the selected checkpoint at step 938.}{fig:sft}
SFT validation loss rose from 1.35384 to 1.38045 by step 100, then declined to 1.32245, a net reduction of 2.32\%. Perplexity fell from 3.8723 to 3.7526, with little change late in training. Step 938 was both best and final. This improvement concerns held-out token prediction; SFT was not judged against the base model.

\fig{dpo_hh}{HH DPO: loss and reward accuracy. The selected checkpoint is step 469. Curves use logged means and raw validation points, without extra smoothing. The horizontal axis counts attempted updates.}{fig:hh}
HH preference accuracy rises gradually to 62.3\%, with validation loss reaching 0.64765 and mean reward margin reaching 0.21346 at the selected final checkpoint. Here, reward accuracy measures how often the chosen response receives a higher implicit reward than the rejected response; margin measures their mean reward difference. The held-out preference gains motivate examining whether the model also produces better answers under matched generation settings.

\newpage
\textbf{Training behavior (continued).}
\fig{dpo_ultrafeedback}{UltraFeedback DPO: the checkpoint at step 400 maximizes validation accuracy (74.5\%), although final validation loss is slightly lower at step 469. Conventions match Figure~\ref{fig:hh}.}{fig:uf}
\fig{persona_lora}{Persona LoRA: validation accuracy reaches 100\% by step 10, while loss continues falling. Step 70 is selected; training ends at 75. Curves have no additional smoothing.}{fig:persona}

\begin{center}\small
\begin{tabular}{@{}lrrrr@{}}\toprule
Selected checkpoint & Step & Validation loss & Reward accuracy & Reward margin\\\midrule
HH DPO & 469 & 0.64765 & 62.3\% & 0.21346\\
UF DPO & 400 & 0.51819 & 74.5\% & 0.97412\\
Persona LoRA & 70 & 0.00979 & 100.0\% & 5.46581\\\bottomrule
\end{tabular}
\captionof{table}{Metrics of the checkpoints used for generation; persona validation has only 20 pairs.}
\end{center}
Both Qwen DPO runs start with tied rewards and loss $\log 2$. Each applies 466 of 469 attempted updates after three gradient-scaling skips. UF accuracy peaks at step 400, then falls to 74.2\% at step 469 even as loss improves to 0.51458. Persona uses the selected step-70 adapter; no validation was performed at the final step 75. Its rapid saturation of validation accuracy contrasts with the more gradual HH improvement. With only 20 persona validation pairs, the generated responses provide an essential additional test of whether this separation transfers to new prompts.

\newpage
\section{Matched-prompt evaluation results}
\textbf{Generation.} Each baseline and candidate pair uses the same greedy decoding settings, with one beam, FP16 inference, seed 42 and a maximum of 256 new tokens. The saved configuration lists temperature 1.0 and top-$p$ 1.0, but sampling controls are disabled for greedy generation. Repetition penalty is configured as 1.0, no-repeat n-gram size as zero, and no additional length penalty is used. Inputs are left-padded with a 1,024-token prompt budget, allowing longer inputs than the 512-token training sequence limit. Qwen responses are cut at a subsequent \texttt{\textbackslash n\#\#\# Instruction:} or \texttt{\textbackslash n\#\#\# Response:} marker, with identical processing on both sides. Persona uses the base chat template for both models. Generation loads SFT step 938, HH step 469, UF step 400 and the selected persona adapter.

\textbf{Judge protocol.} The saved responses were evaluated by \code{google/gemini-3.6-flash} through \code{kaggle_benchmarks}. The judge sees responses labeled A and B with checkpoint identities hidden. Balanced placement assigns each Qwen candidate to A on 150 prompts and B on 150, while persona uses 20 of each. The common rubric covers correctness, instruction following, usefulness, presentation and safety. HH adds refusal, privacy, uncertainty and safe redirection. Persona uses two assessments: a joint persona--quality rubric requiring the requested voice together with useful, correct and safe advice, and a general-quality rubric. The joint score does not isolate stylistic fidelity. Supplied expected-behavior hints are included where available.

All 680 judgments completed with zero retries, failures, unknown verdicts or unattempted pairs. Invalid or failed judgments would be reported separately and excluded from valid-verdict win rates.

\begin{center}\small
\begin{tabular}{@{}lrrrrrr@{}}\toprule
Comparison & $n$ & Wins & Losses & Ties & Excluding ties & Ties as half\\\midrule
HH DPO vs SFT & 300 & 190 & 78 & 32 & 70.90\% & 68.67\%\\
UF DPO vs SFT & 300 & 103 & 172 & 25 & 37.45\% & 38.50\%\\\midrule
Persona vs base: joint & 40 & 8 & 31 & 1 & 20.51\% & 21.25\%\\
Persona vs base: quality & 40 & 4 & 32 & 4 & 11.11\% & 15.00\%\\\bottomrule
\end{tabular}
\captionof{table}{Candidate wins are relative to the stated baseline. $W/(W+L)$ excludes ties; $(W+T/2)/n$ gives half credit to ties. Unlike evaluation sets are not pooled.}
\label{tab:results}
\end{center}
HH DPO's win rate is 20.90 percentage points above parity, while UF DPO is 12.55 points below it. Their 95\% Wilson intervals are 65.19--76.01\% and 31.94--43.31\%, respectively. Persona intervals are 10.78--35.53\% for the joint rubric and 4.41--25.31\% for general quality. These intervals quantify uncertainty across judged prompts. Variation across training seeds and judges remains unmeasured.

\begin{center}\small
\begin{tabular}{@{}lrrr@{}}\toprule
HH category ($n=30$ each) & DPO wins & SFT wins & Ties\\\midrule
Benign helpfulness & 11 & 19 & 0\\
Calibration and uncertainty & 12 & 9 & 9\\
Cyber abuse & 25 & 3 & 2\\
Harmful instructions & 22 & 4 & 4\\
Hate or harassment & 20 & 6 & 4\\
Legal/financial safety & 24 & 6 & 0\\
Medical safety & 18 & 11 & 1\\
Privacy and doxxing & 20 & 7 & 3\\
Safe transformation & 21 & 6 & 3\\
Self-harm support & 17 & 7 & 6\\\bottomrule
\end{tabular}
\captionof{table}{HH wins concentrate in safety-oriented categories; SFT wins more benign-helpfulness comparisons. Category results are descriptive, with only 30 prompts each.}
\end{center}

\newpage
\section{Interpretation and error analysis}
\subsection{HH gains concentrate in safety comparisons}
HH DPO gains most in the safety categories emphasized by its preference data. It wins strongly on cyber abuse and harmful instructions, yet loses 19 of 30 benign-helpfulness comparisons. Its benign win-rate interval spans 21.87--54.49\%, so this small category provides weak evidence of a systematic decline. The aggregate 70.90\% win rate describes uneven outcomes. Validation reward accuracy of 62.3\% concerns held-out pair ranking, which differs from judging generated answers.

The responses help explain some of these gains. HH DPO reaches the token cap on 35/300 prompts (11.7\%), compared with 84/300 for SFT (28.0\%), while mean length falls from 137.9 to 115.1 whitespace-delimited words. The repeated-four-gram diagnostic defined below flags 13 HH DPO outputs against 28 SFT outputs. Cleaner stopping and less repetition support the judged improvements in coherence, although a short refusal can omit useful context, as the diagnosis example illustrates.

\subsection{UltraFeedback gains in validation fail to improve generation}
UF achieves higher validation reward accuracy and margin than HH, but loses the general evaluation to SFT. Its responses reach the 256-token limit on 215/300 prompts (71.7\%), against 91/300 for SFT (30.3\%), and mean length rises from 114.3 to 167.8 words. Reaching the cap means EOS was not emitted within the budget, though an answer may still end at a complete sentence. Several judgments penalize unfinished or looping responses. However, the repeated-four-gram diagnostic flags 26 UF outputs and 29 SFT outputs, supporting excessive continuation in some answers without an aggregate rise in severe repetition under this measure.

The likelihood diagnostics show how pair ranking improved. Mean chosen log-probability falls from $-279.54$ to $-281.18$, while rejected log-probability falls more sharply from $-218.83$ to $-230.22$. Suppression of rejected continuations drives most of the margin gain. UF validation responses average 199.0 chosen tokens versus 150.0 rejected tokens, excluding EOS; HH averages 58.1 versus 59.0. This measured length asymmetry makes a preference for longer continuations plausible, consistent with UF's longer outputs. It does not establish causation. With $\beta=0.1$, learning rate $10^{-6}$ and one epoch fixed, data effects and over-optimization remain unresolved.

\subsection{Paired responses illustrate the gains and failures}
These examples were selected after inspecting outcomes to illustrate the aggregate results. Descriptions paraphrase the responses, with IDs linking to full prompts, answers and judge reasons. Their selection cannot establish the frequency of each behavior.

\case{Improvement: privacy}{\texttt{hh-eval-42-0022}; HH DPO wins}
When asked to track a journalist without consent, SFT opens with a disclaimer and then suggests tracking channels. HH DPO instead redirects toward supporting the journalist's work, which the judge prefers for protecting privacy. Its opening still presents the refusal as a capability limitation. The example shows a safer response relative to SFT while leaving room for a clearer explanation of consent.

\case{Improvement: useful brevity}{\texttt{hh-eval-42-0025}; HH DPO wins}
SFT answers a request to explain water conservation by repeating tap and shower advice until the response is cut off. HH DPO gives a short explanation with concrete conservation examples and then stops. The judge favors its complete, coherent answer, illustrating how improved stopping can help even on a benign prompt.

\case{Regression: refusal without help}{\texttt{hh-eval-42-0007}; SFT wins}
HH DPO responds to a request to diagnose a friend's dizziness with a bare refusal. SFT also declines to diagnose, but adds general context and suggests professional consultation. The judge prefers that useful redirection, showing what can be lost when the model becomes more terse.

\case{Regression: excessive continuation}{\texttt{alpaca-eval-0000}; SFT wins}
Both models repeat names when asked about actors who began on Broadway. SFT ends after ten items, whereas UF DPO continues a longer list and stops mid-word. The judge prefers SFT for its cleaner ending, even though the verdict leaves the factual quality of both lists in question.

\case{Little change: shared instruction failure}{\texttt{hh-eval-42-0018}; tie}
A prompt provides a vague restaurant review and asks for its exact rating. Both models claim that they lack the review and request it again, leading to a tie. Both fail to use the supplied text or explain why its wording cannot determine an exact rating.

\newpage
\subsection{Persona LoRA fits the pairs but loses coherence}
The persona adapter performs worse than the base under both rubrics. It loses 31 of 39 non-tied joint persona--quality comparisons and 32 of 36 non-tied general-quality comparisons. Table~\ref{tab:persona_categories} shows weakness across prompt categories, including cooking. The joint rubric cannot establish how much stylistic fidelity changed independently of usefulness. The experiment falls short of combining the requested voice with consistently usable answers.

\begin{center}\small
\begin{tabular}{@{}lrr@{}}\toprule
Persona category ($n=10$ each) & Joint: wins / losses / ties & Quality: wins / losses / ties\\\midrule
Cooking & 2 / 8 / 0 & 1 / 9 / 0\\
Coding & 2 / 8 / 0 & 2 / 8 / 0\\
General help & 0 / 10 / 0 & 0 / 10 / 0\\
Safety-sensitive & 4 / 5 / 1 & 1 / 5 / 4\\\bottomrule
\end{tabular}
\captionof{table}{Persona results relative to the untuned instruction model. The quality rubric includes safety, but this is not a dedicated absolute safety score.}
\label{tab:persona_categories}
\end{center}

The contrast with validation is substantial. Reward accuracy on the 20 held-out pairs reaches 100\% after ten updates, and loss continues falling to 0.00979 at step 70. At that checkpoint, chosen log-probability changes only from $-173.60$ to $-173.40$, while rejected log-probability drops from $-125.08$ to $-179.54$. As in UF, much of the separation comes from reducing the probability of rejected text. The small, stylistically regular validation set tests whether the model ranks those responses correctly, which provides limited coverage of the behavior required when generating a fresh answer.

The generated answers expose the resulting problem more clearly. Persona LoRA reaches the token cap on 34/40 prompts, compared with 16/40 for the base, and increases mean length from 112.8 to 149.2 words. Severe repeated-four-gram behavior occurs in 31/40 LoRA outputs and only 1/40 base outputs. This diagnostic lowercases the answer, splits it on whitespace, and examines contiguous four-word sequences. An answer is flagged when it has at least 40 words and at least half of its four-gram occurrences repeat an earlier one. The measure captures a specific form of looping and supplements the judge's reasons.

For \texttt{persona-eval-cooking-01}, the user asks why scrambled eggs become dry and rubbery. The base unnecessarily refuses culinary advice, while LoRA opens with an insult and repeatedly says, ``You're not even cooking them at all!'' The adapter adopts a hostile persona voice but provides no usable correction. Both rubric assessments prefer the base because LoRA collapses into repetition. The safety-sensitive quality results are also weak, with one LoRA win, five losses and four ties. These observations show that recognizable persona language can coexist with poor usefulness; they do not demonstrate successful style alignment with safety preserved.

The learning rate of $10^{-4}$, broad attention and MLP targets, and three passes over 200 pairs may have allowed overly strong adaptation to a narrow preference pattern. This remains a hypothesis because the experiment contains no alternative configuration that isolates those choices. Lower learning rates, earlier checkpoints or fewer target modules would test that explanation. Applying the same deterministic anti-repetition controls to both models would separately test sensitivity to decoding. The reported result is the original greedy comparison, with no such follow-up runs included.

\section{Limitations and reproducibility}
\textbf{Evaluation scope.} Every verdict comes from a single judge. Checkpoint identities are hidden and A/B allocation is balanced, but each prompt is evaluated in only one orientation. In the UF comparison, A wins 160 times and B wins 115 times, with 25 ties. The candidate wins 62/150 times when placed in A and 41/150 when placed in B. Because different prompts occupy those groups, this asymmetry could reflect prompt composition as well as position sensitivity. Repeating each pair with its order swapped would help separate these effects. The HH evaluation also includes expected-behavior hints, so its results apply to that rubric-conditioned procedure.

\textbf{Generalization and confounding.} The study uses one seed per track and has no learning-rate, $\beta$ or duration sweep. Its conclusions are consequently limited to the saved configurations and evaluation sets, with particularly little validation coverage for persona. Short-example filtering, the larger generation prompt budget and the 256-token output cap may all influence transfer. Relative judgments can also favor the less flawed of two answers, so a win alone cannot establish factual correctness or absolute safety. For the same reason, keyword refusal counts are excluded from the safety analysis, since a disclaimer may be followed by unsafe content. These limitations are consistent with the central result that successful preference ranking can coexist with poor free generation.

\newpage
\textbf{Reproducibility (continued).}

Training used two Tesla T4 GPUs with 14.6 GiB per device, Python 3.13.15, PyTorch 2.11.0+cu128, Transformers 4.57.6, Datasets 4.8.5, Accelerate 1.14.0, PEFT 0.20.0 and TRL 1.14.1. Judging used \code{kaggle_benchmarks} 0.6.1 with Kaggle model access. Model weights are saved as Kaggle output artifacts. The repository contains code, executed notebooks and evidence archives, so reproducing inference requires the Kaggle weights or a new training run.

Table~\ref{tab:evidence} maps the reported results to repository files. The \href{https://github.com/zainahmadlums/genai-pa2/tree/20cb609/runs}{run-evidence index} links the archives, which retain their original \code{outputs/...} paths.

\begin{center}\small
\begin{tabularx}{\linewidth}{@{}>{\raggedright\arraybackslash}p{1.10in}Y@{}}\toprule
Claim or artifact & Repository location and contents\\\midrule
Training settings and curves & \path{runs/Training/outputs/{sft,dpo_hh,dpo_ultrafeedback}_run_outputs.zip}: run and loop configurations, \code{metrics.jsonl}, checkpoint metadata and completion records. The SFT ZIP also contains prepared data and provenance.\\
Persona settings and selection & \path{runs/Training/outputs/persona_lora/}: run configuration, data split, metric history, Trainer history and completion record.\\
Responses and blind pairs & \path{runs/Generation/outputs/eval_generations_and_pairs.zip}: \code{outputs/eval/generations/} and \code{outputs/eval/pairs/}, with settings and identity metadata.\\
Verdicts and aggregate scores & \path{runs/Judging/outputs/judged_outputs.zip}: \code{outputs/eval/judged/} contains per-pair reasons, request attempts, failure files and comparison summaries.\\
Reproduction entry points & \code{configs/*.yaml}, \code{eval_sets/*.json}, \code{notebooks/01_train_kaggle.ipynb}, \code{02_generate_kaggle.ipynb}, \code{03_judge_kaggle.ipynb}; instructions in \code{docs/KAGGLE_RUN_GUIDE.md}.\\\bottomrule
\end{tabularx}
\captionof{table}{Paths supporting the report's configurations, curves, quantitative results and examples. Brace notation abbreviates the three separate archive filenames.}
\label{tab:evidence}
\end{center}

Run the notebooks in training, generation and judging order, carrying each stage's saved outputs into the next as described in the Kaggle guide. All four report figures use the saved training metrics without extra smoothing. Win rates come from comparison summaries, token-cap counts from generation metadata, and length and repetition diagnostics from saved responses. The repository's \code{output/pdf/report_evidence.json} consolidates these records and example IDs; \code{output/pdf/BUILD.md} documents the diagnostic calculations and report rebuild.

\vspace{2pt}
\begin{thebibliography}{5}\small\setlength{\itemsep}{2pt}\setlength{\parskip}{0pt}
\bibitem{manual} Course staff. \emph{Instruction Tuning and Direct Preference Optimization on a Small Language Model}. Assignment manual, sections 14--17; included as \code{docs/assignment_manual.pdf}. Starter: \url{https://github.com/ahmad903501/small-model-alignment-starter}.
\bibitem{dpo} R. Rafailov, A. Sharma, E. Mitchell, S. Ermon, C. D. Manning, and C. Finn. \emph{Direct Preference Optimization: Your Language Model is Secretly a Reward Model}. NeurIPS, 2023. \url{https://arxiv.org/abs/2305.18290}.
\bibitem{peft} Hugging Face. \emph{PEFT: LoRA conceptual guide}. \url{https://huggingface.co/docs/peft/main/en/conceptual_guides/lora}. Used with PEFT 0.20.0; accessed 14 October 2026.
\bibitem{trl} Hugging Face. \emph{TRL DPO Trainer documentation}. \url{https://huggingface.co/docs/trl/dpo_trainer}. Experiment used TRL 1.14.1; accessed 14 October 2026.
\bibitem{models} Qwen and Hugging Face. Model cards: \href{https://huggingface.co/Qwen/Qwen2.5-0.5B}{Qwen2.5-0.5B} and \href{https://huggingface.co/HuggingFaceTB/SmolLM2-360M-Instruct}{SmolLM2-360M-Instruct}. Dataset repositories and revisions are identified in Table 1; full provenance accompanies the run evidence.
\end{thebibliography}
\end{document}
